\input format.tex

\usepackage{graphicx}
\graphicspath{{cores/}}
\usepackage{setspace}
\usepackage{colortbl,array,booktabs}

\newcolumntype{K}[1]{>{\centering\arraybackslash}p{#1}}

\begin{document}

\vspace*{6mm}

%\color{gray2}
\color{black100}
\noindent\fontsize{20pt}{\baselineskip}\selectfont{肠道有益菌检测报告}

%% 各章节
\setlength{\arrayrulewidth}{1pt}
\fontsize{8pt}{11pt}\selectfont

\vspace*{5mm}

\fontsize{10pt}{13pt}\selectfont
\noindent\begin{tabular}{p{4cm}p{6cm}p{6.2cm}}
\parbox[c]{\hsize}{\vskip7pt 门诊/住院号：测试 \vskip7pt} & 科室：测试 & 送检医生：测试\\
\hline
\parbox[c]{\hsize}{\vskip7pt 姓名：荔湾 \vskip7pt} &
性别：男
& 年龄：35\\
\parbox[c]{\hsize}{\vskip7pt 样本类型：粪便 \vskip7pt} & 样本编号：GI00001 & 送检机构：人和未来\\
\hline
\end{tabular}

\vspace*{6mm}
\begin{spacing}{1.2}
\noindent{\fontsize{15pt}{15pt}\selectfont {有益菌检测结果}}
\end{spacing}
\arrayrulecolor{black100}
\setlength{\arrayrulewidth}{.5pt}

\vspace*{-1mm}
\fontsize{8.8pt}{11pt}\selectfont
\begin{longtable}{|m{4cm}<{\centering}|m{3cm}<{\centering}|m{3cm}<{\centering}|m{5.7cm}<{\centering}|}
\hline
\parbox[c]{\hsize}{\vskip7pt\centering 主要有益菌属 \vskip7pt} &  本次检测值 & 参考范围 & \bf{功能}\\
\hline
\parbox[c]{\hsize}{\vskip7pt\centering 乳酸杆菌属 \vskip7pt} &
0.02
\raisebox{-0.3ex}{\includegraphics[scale=1]{redarrow_down.pdf}}
& ≥ 0.04
& 抑制有害菌的生长，增强人体免疫力 \\
\hline
\parbox[c]{\hsize}{\vskip7pt\centering 双歧杆菌属 \vskip7pt} &
0.01
\raisebox{-0.3ex}{\includegraphics[scale=1]{redarrow_down.pdf}}
& ≥ 0.21
& 抵抗有害菌入侵，促进微量元素吸收 \\
\hline
\parbox[c]{\hsize}{\vskip7pt\centering 阿克曼氏菌属 \vskip7pt} &
未检出
\raisebox{-0.3ex}{\includegraphics[scale=1]{redarrow_down.pdf}}
& ≥ 0.01
& 促进脂肪代谢，预防肥胖 \\
\hline
\parbox[c]{\hsize}{\vskip7pt\centering 柔嫩梭菌属 \vskip7pt} &
13.87
\raisebox{-0.3ex}{\includegraphics[scale=1]{redarrow_down.pdf}}
& ≥ 27.24
& 产丁酸，抑制肠道炎症 \\
\hline
\end{longtable}

%\vspace*{-1cm}
\noindent {结果说明：\raisebox{-0.3ex}{\includegraphics[scale=1]{redarrow_down.pdf}} \ 检测值相比参考人群偏低}

\vspace*{2mm}
\noindent {结果分析：
乳酸杆菌属、双歧杆菌属、阿克曼氏菌属、柔嫩梭菌属结果均异常；\\\\
\vspace*{-1mm}
\hspace*{1.52cm}\begin{minipage}{15.8cm}
\begin{spacing}{1.3}
\fontsize{9pt}{11pt}\selectfont {乳酸杆菌偏低可能导致消化不良、免疫力降低；长期偏低还可引起便秘等病症。}
\fontsize{9pt}{11pt}\selectfont {双歧杆菌含量偏低可能不利于肠道蠕动、营养吸收及免疫调节等；长期偏低还可引起便秘等病症。}
\fontsize{9pt}{11pt}\selectfont {阿克曼氏菌含量偏低，不利于维持能量代谢平衡，还可能影响免疫能力和肠道屏障功能。}
\fontsize{9pt}{11pt}\selectfont {柔嫩梭菌含量偏低，产丁酸减少，不利于抑制肠道炎症和调节免疫功能，长期偏低甚至可能引起克罗恩病、糖尿病等病症。}
\end{spacing}
\end{minipage}
}

%\vspace*{-3mm}
\vspace*{-1mm}
%\fontsize{10pt}{13pt}\selectfont
%\noindent\begin{tabular}{p{2cm}p{5cm}p{2cm}p{5cm}}
%\parbox[c]{\hsize}{\vskip7pt 检测者：} 
% & 审核者：\\
%\parbox[c]{\hsize}{\vskip7pt 检测日期：\vskip7pt} &   & 报告日期：& 2019-12-31\\
%\end{tabular}

\begin{longtable}{m{1.5cm}m{4cm}m{3cm}m{1.5cm}m{2.5cm}}
检测员：&\includegraphics[height=0.9cm,width=3cm]{jianceyuan.png}&&审核员:：&\includegraphics[height=0.9cm,width=3cm]{shenheyuan.png}\\
检测日期：&2019-12-31&&报告日期：& 2019-12-31\\
\end{longtable}
\vspace*{-1mm}
\noindent\fontsize{15pt}{15pt}\selectfont {菌群调理建议} \\

\vspace*{-2mm}
\begin{spacing}{1.3}
\noindent\fontsize{9pt}{11pt}\selectfont
本次结果提示 {\bf 乳酸杆菌属、双歧杆菌属、阿克曼氏菌属、柔嫩梭菌属} 结果异常，建议调节饮食结构、补充益生菌、益生元，促进肠道健康。
\end{spacing}

\setlength{\arrayrulewidth}{.5pt}
\vspace*{1mm}
\begin{center}
\begin{tabular}{p{16.5cm}}
\hspace*{-5mm}
\parbox{1.05\hsize}{
\vspace*{0.5mm}
{\includegraphics[scale=1.0]{xiaoyuandian.pdf}\hspace*{2mm}\fontsize{10pt}{10pt}\selectfont{\bf {膳食建议}}}
}\\

\vspace*{-2mm}
\hspace*{-1mm}
\parbox{1.03\hsize}{
\begin{spacing}{1.5}
\fontsize{9pt}{10pt}\selectfont 建议您在日常饮食中多选择红薯、全麦、洋葱、玉米、蔓越莓等食物，促进柔嫩梭菌属、阿克曼氏菌属、双歧杆菌属等有益菌的生长，增强肠道健康保护力。
\end{spacing}
} \\

\vspace*{1mm}
\hspace*{-5mm}
\parbox{1.05\hsize}{
\vspace*{0.5mm}
{\includegraphics[scale=1.0]{xiaoyuandian.pdf}\hspace*{2mm}\fontsize{10pt}{10pt}\selectfont{\bf {肠道调节建议}}}
}\\

%\vspace*{0.1mm}
\parbox{1.03\hsize}{
\vspace*{1.5mm}
\begin{spacing}{1.5}
\fontsize{9pt}{10pt}\selectfont {\bf {补充益生菌：}} 建议您补充含嗜酸乳杆菌LA-5的乳酸杆菌活菌制剂或者含乳酸杆菌的酸奶、含BB-12的双歧杆菌活菌制剂或者含双歧杆菌的益生菌酸奶等益生菌产品，促进柔嫩梭菌属、阿克曼氏菌属、双歧杆菌属等有益菌的生长，抑制有害菌的异常增殖，调节肠道菌群平衡。 \\
\fontsize{9pt}{10pt}\selectfont {\bf {补充益生元：}} 建议您补充低聚果糖、低聚木糖、富含多酚的蔓越莓提取物、葡萄多酚、菊粉等益生元，促进柔嫩梭菌属、阿克曼氏菌属、双歧杆菌属等有益菌的生长，增加有益物质的产生，有利于增强肠道屏障功能，促进肠道健康。
\end{spacing}
} \\

\end{tabular}
\end{center}


\end{document}

